\LessonHeader{36}{The Dogs That Learn by Listening In}{靠“旁听”学习的狗}{科学新闻 · Science News}{话题\ 动物认知·学说话\enspace|\enspace 新词\ \textbf{12}\enspace|\enspace 新义\ \textbf{6}\enspace|\enspace 短语\ \textbf{14}\enspace|\enspace 全文\ \textbf{831}\ 词}{Lesson 36\enspace·\enspace The Dogs That Learn by Listening In}
\begin{lessonbody}
\PARA{1}Small children \PH{\RV{}{pick}{}\ up}{pick up}{（自然地）学会} a \RV{}{lot}{} of \RVB{}{words}{}{word} just by \RVB{}{listening}{}{listen} to \NWB{}{adults}{}{adult}{/ˈædʌlt/}{n.[C]}{成年人} \RV{}{talk}{.} \RV{}{Sometimes}{} it is the name of a food, and \AG{}{sometimes}{} it is a \RV{}{rude}{} \AG{}{word}{.} \NS{}{Well}{,}{新词性}{}{/wel/}{int.}{嗯；那么} \PH{it\ \RVB{}{turns}{}{turn}\ out\ that}{it turns out that}{原来；结果是} \NW{}{certain}{}{/ˈsɜːtn/}{adj.}{某些} \RVB{}{dogs}{}{dog} can do \RV{}{something}{} very \RV{}{similar}{.}

\PARA{2}\SMK{y}\RVB{}{Scientists}{}{scientist} have \RVB{}{found}{}{find} that \AG{}{dogs}{} with an \RV{}{unusual}{} \NW{}{ability}{,}{/əˈbɪləti/}{n.[sing.]}{能力} the \AG{}{ability}{} to \RV{}{learn}{} the names of \PH{\AGB{}{lots}{}{lot}\ and\ \AGB{}{lots}{}{lot}\ of}{lots and lots of}{非常多的} \RVB{}{objects}{,}{object} can \AG{}{learn}{} those names just by \AG{}{listening}{} to people’s \RVB{}{conversations}{.}{conversation}\EMK{y}{1} 
\YI{1}{科学家发现，有些狗有一种不寻常的本领——能学会许许多多物品的名字。这些狗只靠听人们聊天，就能学会这些名字。}
\RV{}{Nobody}{} has to \RV{}{speak}{} to the \AG{}{dog}{} \DV{}{directly}{.}{direct + -ly → 副词}{adv.}{直接地}

\PARA{3}The \RV{}{team}{} say this is not \RV{}{as}{} \RV{}{simple}{} \AG{}{as}{} it \RVB{}{sounds}{.}{sound} To do it, a \AG{}{dog}{} \RVB{}{needs}{}{need} many \RV{}{social}{} \RVB{}{skills}{.}{skill} It has to \NW{}{identify}{}{/aɪˈdentɪfaɪ/}{vt.}{认出；找出} the \RV{}{important}{} \AG{}{word}{} \RV{}{inside}{} a \RV{}{whole}{} \AG{}{conversation}{.} Then it has to \PH{work\ out}{work out}{弄清楚} what the \AG{}{word}{} is about by \RVB{}{watching}{}{watch} where people look, how they \RV{}{move}{} their \RVB{}{hands}{}{hand} and how their \RVB{}{voices}{}{voice} \AG{}{sound}{.}

\PARA{4}“There are a \AG{}{lot}{} of \RV{}{difficult}{} \AG{}{social}{} things \RVB{}{happening}{}{happen} here, and we don’t know if the \AG{}{dogs}{} are doing it \RV{}{in}{} the \RV{}{same}{} way \AG{}{as}{} the \RVB{}{kids}{,”}{kid} said \PN{}{Dr Shany Dror}{}{/ˌdɒktə ˌʃɑːni ˈdrɔː(r)/}{n.}{沙尼·德罗尔博士} of the \PN{}{University of Veterinary Medicine}{}{/ˌjuːnɪˈvɜːsəti əv ˈvetnri ˈmedsn/}{n.}{维也纳兽医大学} in \PN{}{Vienna}{,}{/viˈenə/}{n.}{维也纳} who \RVB{}{led}{}{lead} the \RV{}{research}{.} “But \PH{on\ the\ \RV{}{surface}{,}}{on the surface}{表面上} the \RV{}{result}{} \RVB{}{seems}{}{seem} to be the \AG{}{same}{.”}

\PARA{5}Writing in the \RV{}{science}{} \RV{}{magazine}{} \PN{}{Science}{,}{/ˈsaɪəns/}{n.}{《科学》（杂志）} the \AG{}{team}{} \RV{}{explain}{} how they made their \RV{}{discovery}{.} They \RVB{}{studied}{}{study} a \RV{}{special}{} \RV{}{group}{} of \RVB{}{animals}{}{animal} that they \RV{}{call}{} \RV{“}{gifted}{} \AG{}{word}{} \RV{}{learner}{”} \AG{}{dogs}{.}

\PARA{6}In the first \RV{}{test}{,} 10 of these \AG{}{dogs}{} were \RV{}{each}{} given two new \RVB{}{toys}{.}{toy} \RV{}{Most}{} of them were \PH{\GL{}{border}{}{/ˈbɔːdə(r)/}{n.[C]}{边境}\ \GL{}{collies}{,}{/ˈkɒli/}{n.[C]}{柯利牧羊犬}}{border collie}{边境牧羊犬} but there was also a \AG{}{border}{} \AG{}{collie}{} \NS{}{mix}{}{新义}{}{/mɪks/}{n.[C]}{混血狗} that was a \PH{\GL{}{rescue}{}{/ˈreskjuː/}{n.[U]}{救助}\ \AG{}{dog}{,}}{rescue dog}{领养的流浪狗} and a \GL{}{Labrador}{.}{/ˈlæbrədɔː(r)/}{n.[C]}{拉布拉多犬} \RV{}{Each}{} \AG{}{toy}{} was \RVB{}{shown}{}{show} to a \AG{}{dog}{} for one \RV{}{minute}{} \RV{}{while}{} \RV{}{someone}{} said its name again and again. Then the \AG{}{dog}{} \PH{was\ \RVB{}{allowed}{}{allow}\ to}{be allowed to do}{被允许做} play with it. \PH{Day\ by\ day,}{day by day}{一天天} this \NW{}{process}{}{/ˈprəʊses/}{n.[C]}{过程} was \RVB{}{repeated}{}{repeat} many times.

\PARA{7}The \RV{}{second}{} \AG{}{test}{} was \RV{}{almost}{} the \AG{}{same}{,} with one big \RV{}{difference}{.} \RV{}{Instead}{} of \AGB{}{showing}{}{show} the new \AG{}{toy}{} to the \AG{}{dog}{,} two \RV{}{family}{} \RVB{}{members}{}{member} \RVB{}{passed}{}{pass} it to each other and \RVB{}{used}{}{use} its name \AG{}{while}{} they \AGB{}{talked}{.}{talk} They made \RV{}{sure}{} they did not play with the \AG{}{dog}{} or \RV{}{even}{} look at it. The \AG{}{dog}{} \RV{}{could}{} only \AG{}{listen}{.}

\PARA{8}\RV{}{Next}{} came the \RV{}{real}{} \AG{}{test}{.} To \RV{}{test}{} \RV{}{whether}{} \AG{}{each}{} \AG{}{dog}{} had \AGB{}{learned}{}{learn} the new names, the \AG{}{team}{} put the two new \AG{}{toys}{} in a \RV{}{room}{} \RV{}{together}{} with nine \RV{}{familiar}{} \AG{}{toys}{.} The \RV{}{owner}{} then asked their \NW{}{pet}{}{/pet/}{n.[C]}{宠物} for one \AG{}{toy}{} by its name, so the \AG{}{dog}{} had to \RV{}{pick}{} it up and \RV{}{bring}{} it back from the \AG{}{room}{.} \AG{}{Each}{} \AG{}{dog}{} \NS{}{completed}{}{新词性}{complete}{/kəmˈpliːt/}{vt.}{完成} 12 \RVB{}{tries}{}{try} for the new \AG{}{toys}{} and 12 for the old ones.

\PARA{9}So how \RV{}{well}{} did they do? \PH{\RV{}{As}{}\ a\ \AG{}{group}{,}}{as a group}{整体来看} the \AG{}{dogs}{} \RVB{}{picked}{}{pick} the new \AG{}{toy}{} \NW{}{correctly}{}{/kəˈrektli/}{adv.}{正确地} about 90\% of the time when they had been told its name \AG{}{directly}{.} When they had only \RVB{}{heard}{}{hear} its name in other people’s \RV{}{talk}{,} they \RV{}{still}{} \PH{got\ it\ \RV{}{right}{}}{get it right}{做对} about 80\% of the time. The \AG{}{scientists}{} say that \AG{}{difference}{} is too small to \RV{}{matter}{.} \SMK{y}In \RV{}{both}{} \RVB{}{cases}{,}{case} the \AG{}{dogs}{} did much better \RV{}{than}{} they \RV{}{would}{} have done just by \RVB{}{guessing}{.}{guess}\EMK{y}{2}
\YI{2}{在两种情况下，狗的成绩都比单靠瞎猜要好得多。}

\PARA{10}“I think the \RV{}{exciting}{} \RV{}{part}{} is what it tells us about these \AG{}{dogs}{’} \AG{}{ability}{} to \RV{}{understand}{} the way we \NW{}{communicate}{,”}{/kəˈmjuːnɪkeɪt/}{vi.}{交流} said Dror.

\PARA{11}A \RV{}{further}{} \RV{}{experiment}{} \RVB{}{suggested}{}{suggest} \AG{}{something}{} \RV{}{even}{} more \RV{}{surprising}{.} \SMK{y}The \AG{}{gifted}{} \AG{}{dogs}{} \AG{}{could}{} \AG{}{learn}{} the name of a new \AG{}{toy}{} \AG{}{even}{} when \AG{}{nobody}{} \NWB{}{mentioned}{}{mention}{/ˈmenʃn/}{vt.}{提到} the name \RV{}{until}{} after the \AG{}{toy}{} had been put where the \AG{}{dog}{} \AG{}{could}{} not see it.\EMK{y}{3} 
\YI{3}{即使直到玩具被放到狗看不见的地方之后才有人提到它的名字，这些有天赋的狗也能学会这个新玩具的名字。}
\RV{}{Young}{} children can do that too.

\PARA{12}The \AG{}{ability}{} to \AG{}{learn}{} names from other people’s \AG{}{talk}{} has been \AG{}{found}{} before in \GL{}{bonobos}{}{/bəˈnəʊbəʊz/}{n.[C]}{倭黑猩猩} that \RVB{}{grew}{}{grow} up in a \RV{}{world}{} \RV{}{full}{} of \RV{}{language}{,} and there is some \RV{}{evidence}{} of it in an \PN{}{African}{}{/ˈæfrɪkən/}{adj.}{非洲的} \RV{}{grey}{} \GL{}{parrot}{.}{/ˈpærət/}{n.[C]}{鹦鹉} But the new \RV{}{study}{} \RVB{}{showed}{}{show} that in the \NW{}{typical}{}{/ˈtɪpɪkl/}{adj.}{典型的；普通的} \AG{}{family}{} \AG{}{dog}{’s} \AG{}{case}{,} the \AG{}{skill}{} is not there.

\PARA{13}\PH{\AG{}{Even}{}\ so,}{even so}{即便如此} the \RVB{}{researchers}{}{researcher} say their \AGB{}{results}{}{result} tell us \AG{}{something}{} \AG{}{important}{:} \RVB{}{humans}{}{human} are not the only \AG{}{animals}{} with the \AG{}{social}{} \AG{}{skills}{} that are \AGB{}{needed}{.}{need} “The \RV{}{fact}{} that this \AG{}{skill}{} also \NWB{}{exists}{}{exist}{/ɪɡˈzɪst/}{vi.}{存在} in an \AG{}{animal}{} that has no \AG{}{language}{} \AGB{}{suggests}{}{suggest} that the \AG{}{skill}{} \RV{}{itself}{} is older \AG{}{than}{} \AG{}{language}{,”} said Dror. “So \AG{}{humans}{} first \RVB{}{developed}{}{develop} an \AG{}{ability}{} to \AG{}{understand}{} \AG{}{difficult}{} \AG{}{social}{} \RVB{}{situations}{,}{situation} and only \RV{}{later}{} \AG{}{used}{} it to \RV{}{build}{} \AG{}{language}{.”}

\PARA{14}We also \RVB{}{heard}{}{hear} from \RV{}{Professor}{} \PN{}{Marilyn Vihman}{}{/ˌmærəlɪn ˈviːmən/}{n.}{玛丽琳·维曼} of the \PN{}{University of California, Berkeley}{,}{/ˌjuːnɪˈvɜːsəti əv ˌkæləˈfɔːniə ˈbɜːkli/}{n.}{加州大学伯克利分校} who did not take part in the work. She said it is \AG{}{still}{} not \NS{}{quite}{}{新义}{}{/kwaɪt/}{adv.}{完全} \RV{}{clear}{} why some \AG{}{dogs}{} are \AG{“}{gifted}{} \AG{}{word}{} \AGB{}{learners}{”.}{learner} But she said their \AG{}{ability}{} to \AG{}{learn}{} \AG{}{words}{} by \PH{\AG{}{listening}{}\ in}{listen in}{旁听} was very \DV{}{impressive}{,}{impress + -ive → 形容词}{adj.}{令人赞叹的} and it \AGB{}{seemed}{}{seem} to \NS{}{match}{}{新义}{}{/mætʃ/}{vt.}{与……相当} what \NWB{}{babies}{}{baby}{/ˈbeɪbi/}{n.[C]}{婴儿} of 18 \RVB{}{months}{}{month} can \RV{}{often}{} do.

\PARA{15}\SMK{y}\AG{“}{Dogs}{} are, \RV{}{as}{} the \RVB{}{writers}{}{writer} \RV{}{point}{} out, very \RV{}{closely}{} \RVB{}{focused}{}{focus} on the \AG{}{humans}{} they \RV{}{depend}{} on to \RV{}{stay}{} \RV{}{alive}{,} so this does not \AG{}{seem}{} \AG{}{surprising}{,”} she said.\EMK{y}{4} 
\YI{4}{她说：“正如作者指出的，狗非常密切地关注它们赖以生存的人类，所以这似乎并不令人意外。”}
But she \RVB{}{added}{}{add} that \AG{}{researchers}{} \RV{}{may}{} have made too many \RVB{}{guesses}{}{guess} \PH{in\ the\ \RV{}{past}{}}{in the past}{在过去} about the \AG{}{social}{} \AG{}{skills}{} that \AG{}{babies}{} \AG{}{need}{} to \AG{}{learn}{} \AG{}{words}{} from other people’s \AG{}{talk}{.} \RV{“}{Perhaps}{} just \AGB{}{hearing}{}{hear} a \AG{}{word}{} again and again is \RV{}{enough}{} for it to \RV{}{stay}{} in the \RV{}{memory}{,} in some \AG{}{cases}{,”} she said.

\PARA{16}So \RV{}{next}{} time you \NS{}{chat}{}{新词性}{}{/tʃæt/}{vi.}{聊天} at home, think about who \RV{}{else}{} \RV{}{might}{} be \AG{}{listening}{.}
\end{lessonbody}
\sourceline{Adapted from: “'Gifted learner dogs' can learn words by eavesdropping, study says”, The Guardian, 2026-01-08.}
\SelfTest{32}{\TW{adult}
\TW{certain}
\TW{ability}
\TW{identify}
\TW{process}
\TW{pet}
\TW{correctly}
\TW{communicate}
\TW{mention}
\TW{typical}
\TW{exist}
\TW{baby}
\TW{well\,\nstag{新词性}}
\TW{mix\,\nstag{新义}}
\TW{complete\,\nstag{新词性}}
\TW{quite\,\nstag{新义}}
\TW{match\,\nstag{新义}}
\TW{chat\,\nstag{新词性}}
\TW{pick up}
\TW{it turns out that}
\TW{lots and lots of}
\TW{work out}
\TW{on the surface}
\TW{border collie}
\TW{rescue dog}
\TW{be allowed to do}
\TW{day by day}
\TW{as a group}
\TW{get it right}
\TW{even so}
\TW{listen in}
\TW{in the past}}
\par\write-1{LESSONEND 36 \thepage}
